\cotefile{57}
\scanpage{20}
Dear Murre,\\

I am glad to hear that you are still willing to \deleted{\unsure{make}}\unsure{give} the talk on unramified functors. Here what I can say to your questions.

1. The theorem about passage to quotient I alluded to is the foll\unsure{owing|hwing}:

\uline{Theorem} Let $f:X\handwritten{\to}Y$ be a morphisme of $S$-preschemes, assume either $X$ and $Y$ \inserted{locally} of finite presentation over $\deleted{Y}\handwritten{S}$, or $Y$ loc noeth and $X$ \inserted{locally} of finite type over $Y\handwritten{\unsure{'}}$. Assume that the equivalence relation $\inserted{R=}X\mathbin{x}_{Y}X$ defined by $f$ is flat over $X$ i.e. $\mathrm{pr}_{1}:X\mathbin{x}_{Y}X\handwritten{\to}X$ is flat. Then the quotient $X/R$ \unsure{exists} in the strongest reasonable sense, \unsure{i.e.} one can factor $f$ into a compositum $X\quad Z\quad Y$, with $\unsure{\mathbb{X}|\mathbb{Y}}\quad Z$ faithfully flat locally of finite presentation, $Z$ locally of finite \deleted{type}\inserted{pres.} over $\deleted{Y}S$ (in fact of finite presentation over $S$ if $X$ is so) and $Z\handwritten{\to}Y$ a monomorphism.

Of course the factorization is unique, andthe theorem can be expressed by saying that the quotient sheaf (for the fpqc topology) $X/R$ is representable. That is in fact how the theorem is proved.

Raynaud has recently made a very nice (and non trivial) application of this theorem, by proving the following: if $S$ is the spectrum of a discrete valuation ring, $G$ a group prescheme of finite type over $S$, $H$ a closed \inserted{and flat} sub-group scheme, such that $G_{t}/H_{t}$ is \deleted{\unsure{affi}} quasi-affine (where $t$ is the generic point of $S$) then $G/H$ is representable as a quasi-affine and flat $S$-scheme, which is even affine if $H$ is invariant (i.e. if $G$ is a flat group scheme of finite type with affine generic fibre, then $G$ is affine). This extends immediately to a b\unsure{a}se which is regular of dim one. Raynaud is now trying to extend his construction to the case when he drops the quasi-affiness assumption, namely to construct still $G/H$
% archive sheet number: 19
